\section{Introduction}

Reinforcement learning has become an increasingly important approach for recommendation systems, as real-world recommendation tasks are inherently sequential and involve long-term interactions between users and platforms. Early recommendation systems were primarily formulated as static prediction problems, where models estimate the relevance of an item to a user based on historical interactions. While supervised learning approaches have achieved notable success in optimizing short-term engagement metrics, they are fundamentally limited in modeling delayed feedback and long-term user satisfaction.

In real-world scenarios, recommendation is inherently a sequential decision-making process. Each recommendation influences future user behavior, preference evolution, and engagement patterns. Optimizing only short-term metrics may lead to suboptimal outcomes such as user fatigue, reduced diversity, or long-term churn. This gap between short-term optimization and long-term value has motivated increasing attention toward reinforcement learning (RL), which naturally models interaction as a Markov Decision Process (MDP) and optimizes cumulative rewards over extended horizons \cite{Sutton1998RL}.

The application of reinforcement learning to recommendation systems has attracted growing interest over the past decade. Early studies explored value-based and policy-based methods to address large action spaces and delayed rewards, including Deep Q-Networks (DQN) and actor–critic frameworks \cite{Zheng2018DRN,Chen2019TopK}. Subsequent research focused on improving state representation, exploration efficiency, and scalability, enabling RL-based methods to better align with industrial-scale recommendation settings \cite{Zhao2018SlateMDP}. More recently, the community has shifted toward deeper forms of exploration and long-term value optimization, supported by high-fidelity simulators and system-level designs \cite{Ie2019DeepExploration}.

In parallel, the rapid advancement of large language models (LLMs) has introduced new possibilities for recommendation systems. By incorporating reasoning, planning, and tool usage capabilities, LLM-powered agents have begun to complement reinforcement learning in sequential recommendation tasks, leading to hybrid architectures that extend beyond traditional model-centric designs \cite{Wang2024RecoMind}. These developments signal a paradigm shift from isolated algorithmic improvements toward holistic, agent-based recommendation systems.

This survey aims to provide a comprehensive and structured overview of reinforcement learning-based recommendation systems. Specifically, we make the following contributions:
\begin{itemize}
    \item We formalize recommendation systems from a reinforcement learning perspective and summarize common problem formulations.
    \item We propose a taxonomy that categorizes RL-based recommendation methods according to learning paradigms, state representation strategies, and exploration mechanisms.
    \item We review representative works to highlight key technical innovations and practical lessons.
    \item We analyze emerging trends, including deep exploration and the integration of RL with large language models, and discuss open challenges and future directions.
\end{itemize}

By synthesizing existing research through a unified lens, this survey seeks to facilitate a deeper understanding of how reinforcement learning can address the core challenges of modern recommendation systems.
\label{sec:intro}